\section{Results}
\label{sec:results}

\subsection{No Single Observer Dominates Across Systems and Metrics}
\label{sec:results:overview}

Table~\ref{tab:benchmark_500} summarizes the 500-patient benchmark.

On the fold system, Kalman-LSTM-Spec achieves higher early-warning AUC
(\(0.903 \pm 0.059\) vs.\ \(0.827 \pm 0.235\)) with lower cross-seed variance.
Kalman-BCE produces higher mean detection time (\(109.5 \pm 28.6\) vs.\
\(72.7 \pm 11.0\) steps), indicating earlier threshold crossings, and a
slightly lower false-positive rate (\(0.330\) vs.\ \(0.379\)).

On the Hopf system, both methods maintain low false-positive rates
(\(\approx 0.05\)). LSTM-Spec again has higher EW-AUC (\(0.823 \pm 0.243\)
vs.\ \(0.617 \pm 0.177\)). BCE has higher mean detection time (\(60.3\) vs.\
\(49.5\) steps) but with larger variance (\(\pm 26.5\) vs.\ \(\pm 3.8\)).

On the logistic system, both methods reach perfect EW-AUC (\(1.000\)) and
similar detection times (\(\approx 62\) steps). The only separation is
false-positive rate: BCE is lower (\(0.201 \pm 0.143\) vs.\
\(0.268 \pm 0.098\)).

No method dominates all three metrics on any system. LSTM-Spec tends to
discriminate better (higher EW-AUC on fold and Hopf), while BCE is preferable
when false-positive control is the priority or when earlier raw detection time
matters.

\subsection{Patient Count Changes the Method Ranking}
\label{sec:results:sweep}

Table~\ref{tab:benchmark_sweep} and Figure~\ref{fig:patient_sweep} report
performance across patient counts from 100 to 500.

On the fold system, LSTM-Spec detection time decreases from \(108\) to \(73\)
steps as patient count grows, while BCE detection time stays near
\(110\)--\(121\). The EW-AUC gap widens in the same direction: both methods
start near \(0.87\)--\(0.89\) at 100~patients, but by 500~patients LSTM-Spec
reaches \(0.903\) while BCE is at \(0.827\). LSTM-Spec benefits more from
additional training data on this system.

On the Hopf system, LSTM-Spec EW-AUC increases from \(0.652\) to \(0.823\)
across the sweep, compared with a flat range of \(0.60\)--\(0.64\) for BCE.
However, LSTM-Spec detection time is undefined at 200~patients and is reported
from a single seed at 300~patients. The threshold (\(0.500\)) exceeds the
model's pre-bifurcation scores at these patient counts, even though the scores
still discriminate in the EW-AUC sense. By 400--500~patients, threshold
crossings become consistent.

On the logistic system, both methods reach EW-AUC~\(= 1.0\) by 200~patients
and maintain stable detection times near \(62\)--\(66\) steps. False-positive
rate increases with patient count for both methods, with LSTM-Spec consistently
higher.

\subsection{Score Trajectories Reveal Regime-Specific Behavior}
\label{sec:results:trajectories}

Figure~\ref{fig:trajectory_panels} shows representative score trajectories at
500~patients.

On the fold system, the BCE score rises gradually and crosses threshold early in
the trajectory. The LSTM-Spec score holds low longer and then rises sharply near
the bifurcation. This matches the aggregate pattern: BCE has higher mean
detection time (earlier crossings) but wider variance, while LSTM-Spec has a
tighter, later crossing with better signal--null separation.

On the Hopf system, the BCE score is nearly flat. The LSTM-Spec score shows more
structure near the bifurcation but, in the displayed seed, does not cross
threshold before the transition.

On the logistic system, the BCE score remains below threshold for the displayed
seed. The LSTM-Spec score crosses threshold early and drops after the
bifurcation point.

\subsection{Training Curves and Optimization Behavior}
\label{sec:results:training}

Figure~\ref{fig:training_curves} plots training and validation loss for both
variants at 500~patients. Both methods converge smoothly with minimal
train--validation gap. BCE plateaus around epoch~20--30, while LSTM-Spec
continues to decrease through the full 50~epochs. LSTM-Spec reaches a lower
final loss on all three systems.

The lower training loss for LSTM-Spec is consistent with its higher EW-AUC on
fold and Hopf: the model fits the Gaussian ramp target more tightly. On the
logistic system, both methods reach perfect EW-AUC despite a large loss
difference, indicating that the logistic transition is detectable even with a
coarser fit to the training target.

\subsection{Summary of Main Findings}
\label{sec:results:summary}

\begin{itemize}[leftmargin=*]
  \item LSTM-Spec achieves higher EW-AUC than BCE on fold and Hopf at
        500~patients. Both reach perfect EW-AUC on logistic.
  \item BCE produces earlier detections (higher DT) on fold and comparable DT
        elsewhere, but with larger cross-seed variance.
  \item BCE has lower false-positive rate on fold and logistic. Both methods
        have low FPR on Hopf at high patient counts.
  \item LSTM-Spec benefits more from increasing patient count on fold and Hopf.
        Logistic performance saturates by 200~patients for both methods.
  \item Score trajectories differ in shape: BCE rises gradually, LSTM-Spec
        produces sharper transitions closer to the bifurcation.
\end{itemize}
